% ====================================================================
% The two-factor law: why the router (sometimes) matters.
% Replaces the retracted cut-based argument (which predicted the wrong
% sign of the tube/sp-divergence relation and was removed). Everything
% here is measured: activity audit, interventional capacity sweeps,
% relief sweep, partial correlations. No claim without an experiment.
% ====================================================================
\section{Why the Router (Sometimes) Matters: A Two-Factor Law}
\label{sec:twofactor}

Section~\ref{sec:attractor} established that congestion-aware routers settle
into nearly one load distribution. This section asks when the residual
differences between them matter, and answers with a law in two factors---one
of state, one of structure---each demonstrated by intervention or by partial
correlation rather than by cross-sectional association alone.

\subsection{The microfoundation: physical terms are congestion-gated}
\label{sec:gating}

The added terms of the physical cores are not always active. The Newton scar
requires drops to exist at all; the Archimedean buoyancy requires edge density
to exceed its flotation threshold. On topologies with negligible loss, both
terms are mathematically zero and those cores \emph{are} the bare substrate:
on our zero-drop scale-free graph, both relocate $0.00\%$ of load mass
relative to it. Only the Pascalian diffusion, which reads neighbour load
rather than failures, stays active everywhere ($3$--$10\%$ of mass moved,
uncorrelated with drop level). An activity audit across the bench shows the
scar moving $0$--$16\%$ of load mass and the buoyancy $0$--$3.3\%$, both
tracking congestion. Sweeping the traffic-pressure axis confirms the same conditioning at the level of the law: as link capacity is raised until no router drops a packet, the advantage of congestion-aware routing over the blind baselines vanishes to zero---with no losses to redistribute, the law has nothing to predict---so the law carries signal only in the regime where losses exist to be moved. Two consequences follow. The convergence of
Section~\ref{sec:attractor} must be read with this conditioning: where the
gated terms are inactive, agreement between those cores is degeneracy, not
evidence; the meaningful convergence is the one observed where all mechanisms
are demonstrably active (the high-congestion topologies), and it holds there
too. And mechanism divergence should rise with congestion---which we can test
causally.

\subsection{Factor one, by intervention: congestion drives divergence}
\label{sec:causal}

Cross-topology correlations confound everything topologies differ in. We
therefore held each topology fixed and turned a single knob: every edge
capacity scaled by a factor $f$, the demand and seeds unchanged. Across four
topologies and six pressure levels, inter-core divergence tracks the drop rate
within every graph (Spearman $+0.77$ to $+1.00$). The endpoint reversals are
the sharpest statement: the scale-free graph on which the cores agree almost
perfectly diverges by a factor of $21$ once squeezed ($f=0.45$), and G\'EANT,
the bench's most divergent graph at native capacity, converges by a factor of
$8$ once relieved ($f=1.5$). Same graph, same demand, opposite outcomes:
congestion is the cause, not a correlate, of mechanism divergence. The sweep
also reveals a saturation regime: past heavy loss ($f\le0.6$, drop rates above
$\sim35\%$) divergence plateaus or falls again---when every path drowns, no
mechanism can act, and the distributions reconverge in collapse. Mechanism
identity matters in the \emph{contested} band between free flow and gridlock.

\subsection{Factor two: structure sets the ceiling}
\label{sec:ceiling}

The real Abilene backbone is the stress case that exposes the second factor.
It is by far the most congested graph on the bench (drop rate $0.22$ at native
capacity), yet its inter-core divergence is unremarkable ($0.017$)---it
single-handedly breaks any one-factor congestion law across topologies. The
reason is structural: with eleven nodes and fourteen edges it contains four
cycles, and divergence between routing mechanisms can only be expressed
through alternative paths that exist. A capacity sweep on Abilene itself
confirms both factors at once: relieving it traces the same inverted-U as
every other graph (divergence rises from $0.001$ in free flow to a peak of
$0.030$ in the contested band, then falls in saturation), but the peak is a
third of what structurally richer graphs reach under comparable pressure.
Congestion dictates the \emph{shape} of the divergence curve---where on it a
network sits; structure dictates its \emph{amplitude}---how high the curve can
go.

\subsection{The dissociation: two predictors, two distinct roles}
\label{sec:dissociation}

The two factors separate cleanly in partial correlation, on opposite
questions. For \emph{mechanism divergence}, congestion is the driver
($r=+0.92$ controlling for tube/sp) and tube/sp retains nothing ($r=-0.10$
controlling for congestion). For \emph{performance gains} over blind
baselines, the roles invert: tube/sp retains a significant independent link
($r=+0.59$ controlling for congestion) while congestion retains none
($r=-0.23$). One variable of state, one of structure, each owning one
question: congestion determines \emph{how much the choice of mechanism
matters}; structure determines \emph{how much any congestion-aware rerouting
can gain}. Both are measurable before deployment---one from a traffic
estimate, one from the graph alone.

\subsection{The law, stated}

\emph{The load distribution of a graph under demand is an attractor that any
congestion-aware rule approaches. The identity of the rule matters only in the
contested regime between free flow and saturation, in proportion to the
congestion it must manage (factor of state, causal), and only up to the
diversity of paths the structure provides (factor of structure, the ceiling).}
This is an empirical law with a mechanism (term gating), interventional
support (capacity sweeps in both directions), and a stated boundary; a formal
derivation of the attractor and of the ceiling's dependence on cycle structure
remains open, and we offer it as the target for theoretical treatment. One candidate has already been eliminated: a cut-based account of the attractor was formulated and tested during this work, predicted the wrong sign of the structure--divergence relation, and was retracted---the open problem is therefore less trivial than it may appear.
